\begin{figure}[tbp]
\centering
\ViewHeading{(a) The supplied datum and its quotient seed}
\begin{tikzpicture}
\node[vinput] (seed) at (-4.8,0) {$S$};
\node[vbox] (cycles) at (0,0) {$\mathsf Z$};
\node[vresult] (quotient) at (4.8,0) {$H=\mathsf Z/\mathsf B$};
\draw[varrow] (seed)--(cycles) node[midway,above] {$\widetilde\delta$};
\draw[varrow] (cycles)--(quotient) node[midway,above] {$\pi$};
\draw[varrow] (seed.south) to[bend right=22] node[midway,below] {$\ob=\pi\widetilde\delta$} (quotient.south);
\node[font=\footnotesize] at (4.8,-1.55) {$V=\operatorname{im}\ob\subseteq H$};
\end{tikzpicture}
\ViewHeading{(b) Boundary preservation is what makes descent well defined}
\begin{tikzpicture}
\node[vbox] (upperleft) at (-3.3,0) {$\mathsf Z$};
\node[vbox] (upperright) at (3.3,0) {$\mathsf Z$};
\node[vresult] (lowerleft) at (-3.3,-1.7) {$H$};
\node[vresult] (lowerright) at (3.3,-1.7) {$H$};
\draw[varrow] (upperleft)--(upperright) node[midway,above] {$G_j$};
\draw[varrow] (lowerleft)--(lowerright) node[midway,below] {$\overline G_j$};
\draw[varrow] (upperleft)--(lowerleft) node[midway,left] {$\pi$};
\draw[varrow] (upperright)--(lowerright) node[midway,right] {$\pi$};
\node[vinput,text width=11cm] at (0,-3.2)
 {$z'=z+b$, $b\in\mathsf B$\quad$\Longrightarrow$\quad
  $G_jz'-G_jz=G_jb\in\mathsf B$\\
  hence $[G_jz']=[G_jz]$};
\end{tikzpicture}
\caption{The response starts with a specified seed map and an indexed
tuple of marks. In this figure $H=\Hex$. Relative exactness can supply
the quotient seed, but does not select the endomorphisms. The square
commutes precisely because the displayed marks preserve the boundary
subspace.}
\label{fig:response-datum}
\end{figure}